\section{Results}
\label{sec:results}

The results in this section retain their August 2026 experimental scope.
The September public-configuration audit and CT preview are reported
separately in Section~\ref{sec:public-release}; the coverage measures and
input lineages should not be combined.

\plain{In brief: we report the stages separately --- sheet coverage, reading order, winding energy, metric distortion, collision state, and raster quality --- because they fail in different ways, and a single aggregate would hide the one failure that matters. The experiments that did not work are reported alongside those that did.}

A single aggregate score is the wrong instrument for this problem, for the reason Section~\ref{sec:gp} makes concrete: the decisive failure --- a surface that has stepped onto the neighbouring turn --- is invisible to every geometric quality measure, so any score computed from geometry can be improved by committing it. Each stage below is therefore measured against a quantity it could have failed, and where a stage is judged against CT, the fit being judged never saw that CT.

\subsection{Sheet Coverage and Reading Order}

The sheet books of Table~\ref{tab:sheet-book} carry more than 70\% of evidenced material on every rung. Count and material diverge sharply at full scale: the $21^3$ book accepts 6,188 of 6,500 lifted sheets, yet its largest single sheet carries only 0.15\% of the material and the 70\% prefix requires 5,493 tracks. That fragmentation is a property of the evidence, and the held-out CT check catches both local bad interpolation and large prediction phantoms.

Winding assembly scales more favourably than the geometry count. The $21^3$ main component orders 15,457 tracks over 112 winds at 5.0\% residual conflict weight, assigns 99.3\% of accepted material, and suspends the order claims of 174 tracks; the smaller $4\times5\times5$ and $10\times10\times10$ rungs need none. The interpretation is geometric: radial order survives fragmentation that geometry does not, because it is measured from separation across a local pitch rather than from connectivity. The remaining difficulty is longitudinal, not radial.

\subsection{Probabilistic Winding Results}

The generalized winding field supplies an independent, noisy absolute cue to the discrete relation graph. In a full-field probe of 313,615 vertices, the two-sided field exhibits the clean-sheet unit jump at 310,686 vertices (99.1\%); the remaining 0.9\% are discarded principally by local orientation cancellation. This measures local field behaviour, not an absolute turn label. Absolute gauge is still fixed once per relation island, and the field is downweighted or discarded where its within-island calibration is weak.

On a fresh $4\times5\times5$ PHerc0139 ball-pivoted pile (882,219 vertices, 1,711,155 faces, and 783 disk components), the standalone winding path contracts 783 patches into 25 continuation components and four relation islands. Of 2,620 measured relations, 1,698 enter the MRF. Two recentered graph-cut rounds change four component labels, reduce the measured energy from 221.369 to 200.486, and leave continuation and order satisfaction at 0.99971 and 0.99816, respectively; the complete winding-only run takes 10.8 seconds. Mean conditional label confidence is 0.99992, but we report it only as local sharpness with neighbouring MAP labels fixed, not as an exact global marginal.

The outer conflict-penalty experiments establish the present boundary. An unguarded 24-round run is rejected because its winding span expands from 14.13 to 32.41 turns --- the penalties bought local agreement by inventing turns. A broad sweep finds 21,420 cells in which spatially remote geometry lands in the same developed position; winner- and context-lock variants reduce this to 11,099/10,214 only by breaking the relations the labels were built from, dropping continuation/order satisfaction to 0.95789/0.66605 and 0.95580/0.80329. Restricting attention to cells with at least four mutually remote claimants leaves two cells and four losing claims. The one-bin variant clears both and preserves winding span, but creates 42 residual relation conflicts and raises MRF energy from 200.486 to 1,119.330; requiring the loss in two cells makes no exclusion at all and preserves the base labels bit-for-bit. The base winding result stands; no penalized configuration improves duplicate claims without damaging the relations.

\subsection{Metric Projection Across Frozen Rungs}

Table~\ref{tab:metric-rungs} reports the common fixed-topology projection engine on three frozen inputs. It is explicitly \emph{not} an end-to-end scale table: $4\times5\times5$ carries authoritative phase from the modern fit, $10\times10\times10$ uses an older solid input, and $4\times21\times21$ uses a carry-registered ribbon. What the comparison does test is whether one metric projection preserves each input's established order without a new gauge search --- that is, whether the projection is a function of its input rather than of its tuning.

\begin{table*}[t]
  \caption{Fixed-topology metric projection on three frozen upstream states. Fill is reported only where the source run defines a comparable raster; $4\times5\times5$ instead reports the composited painted population. These are three different inputs through one engine, not one pipeline at three scales.}
  \label{tab:metric-rungs}
  \centering
  \small
  \begin{tabular}{@{}lrrrrrr@{}}
    \toprule
    Grid & Vertices & Faces & Output coverage & Length RMS & Link $|\Delta U|$ med./p95 & Time \\
    \midrule
    $4\times5\times5$ & 2.040M & 4.033M & 6.82M painted px & 0.00452 & 0.300 / 1.15 & 9.6 s \\
    $10\times10\times10$ & 1.512M & 3.012M & 60.8\% fill & 0.0216 & 6.02 / 48.3 & 6.8 s \\
    $4\times21\times21$ & 6.731M & 12.946M & 68.1\% fill & 0.00415 & 0.205 / 1.27 & 30.8 s \\
    \bottomrule
  \end{tabular}
\end{table*}

The $4\times5\times5$ map has $|du/ds-1|$ at numerical precision and transfers every vertex and face. Phase placement moves 66 detached components from an atlas-like packed input into physical order, and secondary layers recover every marked missing region while keeping 1.25 million duplicated-ply pixels out of the primary view (Fig.~\ref{fig:metric-sheet}). The $4\times21\times21$ carried frame is already close to ideal and remains so. The $10\times10\times10$ frame is substantially compressed and rough --- the common metric map expands its mean rate from 0.777, and residual cross-row jitter and source-fused continuation remain --- but even there it improves the retired solver's link $p95$ from 89.6 to 48.3 and alignment RMS from 290 to 150 \emph{without changing carried order}, which is the property being tested.

The unified pre-weld path introduces the robust consistency system because its independent cube observations contain cracks the frozen projection inputs do not. At $4\times5\times5$, 2.04 million unknowns and 12.2 million nonzeros factor in 8--11 seconds; the crack rows reduce local slips from $\pm20$--25 to $\pm1.3$ pixels, and median texture seam shift falls to 3 pixels. The canonical $4\times5\times5$ sequence completes end to end. A separate $21\times21\times21$ campaign reaches the maintained pre-fit concatenation boundary with 9,110 unique meshes and intentionally stops there; the fixed-topology comparison numbers above remain attached to the experiments that produced them.

A preliminary post-projection experiment asks whether continuous distortion can improve without reopening the discrete winding decision. On the boundary-complete 882,219-vertex, 1,711,155-face pile, eight UV-only symmetric-Dirichlet sweeps with a two-voxel displacement cap take 8.1 seconds. Area-weighted mean anisotropy falls from 1.761554 to 1.628350 and symmetric-Dirichlet energy from 63.28M to 56.95M; the independent primary-sheet raster improves condition-number median/p95/p99 from 1.408/2.661/6.014 to 1.320/2.413/5.278. Mean displacement is 0.190 voxels, no face changes its input orientation, and the 24,340 faces already over the stretch bound are left untouched. This is an order-preserving result on one rung, not a pipeline default or a scaling claim.

\subsection{Attachment, Distortion, and Output Quality}

Figure~\ref{fig:recto-fixed-point} and Section~\ref{sec:attach} give the mature $4\times5\times5$ attachment result: V83 carries 8.168 million faces with zero coverage fault, 3.162\% dark pixels, Sander $p95/p99$ 1.190/1.380, and 428 samples at $\ge4\times$. Against the fixed progressive baseline, the V46/V83 lineage removes multiple coverage entirely and has roughly 320 times fewer $\ge4\times$ samples. The loop terminates by rejecting its own final proposal, every measured quality tail having worsened, which locates the fixed point by measurement rather than by a visually chosen iteration count.

\paragraph{Provenance.} One piece of bookkeeping earns its place among the results, because it changes what a coverage number means. Every vertex and pixel carries a label separating directly supported geometry from harmonic or cylindrical interpolation, additive prediction bridges, and vertices moved by CT attachment. The classes behave differently, and averaging them hides the difference. On $10\times10\times10$, restricting the output to fitted material plus interpolation that survives its final CT and distortion state gives 70.1\% coverage at 3.59\% dark pixels, Sander $p95$ 1.45, and a 0.03\% $\ge4\times$ tail; the full rectangle --- which contains continuation across wraps the source prediction had already fused --- reports a near-identical 70.6\% coverage at 4.58\% dark, Sander $p95$ 2.73, and a 2.53\% tail. Coverage is at parity; every quality measure is not, and only the labels distinguish the two. On $4\times5\times5$ the same restriction yields 83.2\% coverage, 3.02\% dark, and no $\ge4\times$ samples.

\subsection{Measured Dead Ends and Remaining Failure Modes}
\label{sec:deadends}

Five negative results narrowed the method, and each one fixed a design decision described earlier in the paper. What follows is in each case an account of what the objective actually optimizes, because in every one of these the implementation was correct and the objective was answering a different question from the one intended.

\begin{itemize}
  \setlength{\itemsep}{3pt}
  \item \textbf{Winding-conflict penalties (Section~\ref{sec:prob-winding}).} A near-hard unary forbidding label $k$ for component $c$ has a cheapest escape, and it is the wrong one. The pairwise factors are local, so translating an entire weakly connected component several turns away in $k$ costs almost nothing, while satisfying the forbidden assignment in place costs the relations that hold it. The penalty's least-energy solution is therefore to invent turns --- which is exactly the failure it was added to prevent, and precisely what the unguarded run does when the span goes from 14.13 to 32.41. Restricting which cells may issue a penalty bounds the damage without changing its direction: broad winner and context locks halve the conflicting cells by breaking the relations the labels rest on, the one-bin variant clears its two four-claimant cells at more than fivefold the energy, and the two-bin variant issues no penalty at all. The lesson is about the relative strength of two measurements, not about tuning: a coincidence in developed position is weaker evidence than a measured radial relation, and an objective that lets the first overrule the second will spend the second to satisfy it.

  \item \textbf{The geometric slab splitter.} It carves 11.4 million prediction pixels (4.1\% of foreground) along visually convincing medial ridges in 2.8 seconds and separates 0 of 246 measured fused runs. The reason is a mismatch of scales between two operators: a cut in the probability volume is one voxel wide, while reconstruction spans gaps up to the pivot radius $\rho=1.2$ and the MLS collapse first pulls points from both sides of the cut toward a common surface. A separation the mesher cannot see is not a separation. Cutting in the voxel domain can only work if the cut exceeds the reconstruction's own reach, and a medial ridge never does.

  \item \textbf{The tabu winding atlas (Appendix~\ref{app:tabu}).} It reduces 373,536 exact UV overlaps to 1,874 --- a 99.5\% reduction --- while tearing 41 usable weld relations. It remains a reproducible baseline, and its failure is the reason the current system does not let overlap minimization choose winding identity: overlap count is a geometric objective, and by Section~\ref{sec:gp} a geometric objective cannot see the distinction it was being asked to decide. At $10\times10\times10$ the search is net-negative under every configuration tested, and on both tested rungs the no-search control plus $u$-packing reaches higher coverage than any search configuration (71.7\% and 70.1\%): the deficit the search was introduced to close was packing, not placement.

  \item \textbf{The full elastic shell (Section~\ref{sec:untangle}).} It reduces 147,728 residual contacts to 4,073 after snap and then stops. Each contact contributes a unilateral clearance constraint on a scalar radial field, and on a tightly wound coil the constraints are not independent: opening one penetration presses the sheet against the next turn, whose constraint then forbids the motion. Past a certain contact density the feasible set is small enough that no admissible displacement remains, so the residue is a property of the constraint system rather than of the number of rounds. Contact pruning and the newer pre-weld input are the next controlled variables.

  \item \textbf{Generated continuation at $10\times10\times10$.} The fitted region is high quality, but continuation crosses wraps that were already fused in the source prediction. A parameterizer receives a surface, not a history: where two turns arrived already joined, there is no residual in any energy it minimizes that distinguishes that join from a real one, and a smooth continuation across it is the correct answer to the question it was asked. This is the direct motivation for the division of labour with the companion pipeline: the distinction is representable in the volume, where the two sheets still have different radii, and is gone once they are one component. No amount of care downstream recovers it, which is why the prediction volume is completed before it reaches us rather than repaired after meshing.
\end{itemize}

One pattern runs through all six. In each case an operator was asked for a decision whose evidence had already been destroyed by an earlier representation change --- reconstruction closing a one-voxel cut, an overlap count standing in for a measured relation, a parameterization receiving two turns already fused. The gains came from moving each decision to the domain where the quantity it turns on is still present: winding phase in the mesh, radial order in the relation graph, arc length in the ribbon. Strengthening the downstream optimizer never helped, because the information it needed was not in its input.
